%%%%%%%%%%%%%%%%%
% This is an sample CV template created using altacv.cls
% (v1.7.4, 30 Jul 2025) written by LianTze Lim (liantze@gmail.com), based on the
% CV created by BusinessInsider at http://www.businessinsider.my/a-sample-resume-for-marissa-mayer-2016-7/?r=US&IR=T
%
%% It may be distributed and/or modified under the
%% conditions of the LaTeX Project Public License, either version 1.3
%% of this license or (at your option) any later version.
%% The latest version of this license is in
%%    http://www.latex-project.org/lppl.txt
%% and version 1.3 or later is part of all distributions of LaTeX
%% version 2003/12/01 or later.
%%%%%%%%%%%%%%%%

%% Use the "normalphoto" option if you want a normal photo instead of cropped to a circle
% \documentclass[10pt,a4paper,withhypeper,normalphoto]{altacv}

\documentclass[10pt,a4paper,withhyper]{altacv}
%% AltaCV uses the fontawesome5 and simpleicons packages.
%% See http://texdoc.net/pkg/fontawesome5 and http://texdoc.net/pkg/simpleicons for full list of symbols.

% Change the page layout if you need to
\geometry{left=1.25cm,right=1.25cm,top=1.5cm,bottom=1.5cm,columnsep=1.2cm}

% The paracol package lets you typeset columns of text in parallel
\usepackage{paracol}


% Change the font if you want to, depending on whether
% you're using pdflatex or xelatex/lualatex
% WHEN COMPILING WITH XELATEX PLEASE USE
% xelatex -shell-escape -output-driver="xdvipdfmx -z 0" mmayer.tex
\iftutex
  % If using xelatex or lualatex:
  \setmainfont{Lato}
\else
  % If using pdflatex:
  \usepackage[default]{lato}
\fi

% Change the colours if you want to
\definecolor{VividPurple}{HTML}{3E0097}
\definecolor{SlateGrey}{HTML}{2E2E2E}
\definecolor{LightGrey}{HTML}{666666}
% \colorlet{name}{black}
% \colorlet{tagline}{PastelRed}
\colorlet{heading}{VividPurple}
\colorlet{headingrule}{VividPurple}
% \colorlet{subheading}{PastelRed}
\colorlet{accent}{VividPurple}
\colorlet{emphasis}{SlateGrey}
\colorlet{body}{LightGrey}

% Change some fonts, if necessary
% \renewcommand{\namefont}{\Huge\rmfamily\bfseries}
% \renewcommand{\personalinfofont}{\footnotesize}
% \renewcommand{\cvsectionfont}{\LARGE\rmfamily\bfseries}
% \renewcommand{\cvsubsectionfont}{\large\bfseries}

% Change the bullets for itemize and rating marker
% for \cvskill if you want to
\renewcommand{\cvItemMarker}{{\small\textbullet}}
\renewcommand{\cvRatingMarker}{\faCircle}
% ...and the markers for the date/location for \cvevent
% \renewcommand{\cvDateMarker}{\faCalendar*[regular]}
% \renewcommand{\cvLocationMarker}{\faMapMarker*}


% If your CV/résumé is in a language other than English,
% then you probably want to change these so that when you
% copy-paste from the PDF or run pdftotext, the location
% and date marker icons for \cvevent will paste as correct
% translations. For example Spanish:
% \renewcommand{\locationname}{Ubicación}
% \renewcommand{\datename}{Fecha}


%% Use (and optionally edit if necessary) this .tex if you
%% want to use an author-year reference style like APA(6)
%% for your publication list
% \input{pubs-authoryea}

%% Use (and optionally edit if necessary) this .tex if you
%% want an originally numerical reference style like IEEE
%% for your publication list
\input{pubs-num}

%% sample.bib contains your publications
\addbibresource{sample.bib}

\newcommand{\cvpub}[3]{%
  {\normalsize\color{accent} #1}\par
  \vspace{2pt}
  {\small #2}\par
  \vspace{2pt}
  {\small\itshape #3}\par
  \vspace{6pt}
}

\begin{document}
\name{Sieun Park}
\tagline{Mobile Systems Researcher | AI Privacy · Accountable Mobile AI}
% Cropped to square from https://en.wikipedia.org/wiki/Marissa_Mayer#/media/File:Marissa_Mayer_May_2014_(cropped).jpg, CC-BY 2.0
%% You can add multiple photos on the left or right
\photoR{3cm}{id_photo.jpg}
% \photoL{2cm}{Yacht_High,Suitcase_High}
\personalinfo{%
  % Not all of these are required!
  \email{sieun.park@hcs.snu.ac.kr}
  \phone{+82 10-5962-2446}
  \location{Seoul, South Korea}
  \linkedin{sieunpark-at-snu}
  % \twitter{@marissamayer}
  \homepage{cherry-0.github.io/sieunpark-at-snu/}
%   \github{github.com/mmayer} % I'm just making this up though.
%   \orcid{0000-0000-0000-0000} % Obviously making this up too.
  %% You can add your own arbitrary detail with
  %% \printinfo{symbol}{detail}[optional hyperlink prefix]
  %% \printinfo{\faPaw}{Hey ho!}
  %% Or you can declare your own field with
  %% \NewInfoFiled{fieldname}{symbol}[optional hyperlink prefix] and use it:
  % \NewInfoField{gitlab}{\faGitlab}[https://gitlab.com/]
  % \gitlab{your_id}
	%%
  %% For services and platforms like Mastodon where there isn't a
  %% straightforward relation between the user ID/nickname and the hyperlink,
  %% you can use \printinfo directly e.g.
  % \printinfo{\faMastodon}{@username@instace}[https://instance.url/@username]
  %% But if you absolutely want to create new dedicated info fields for
  %% such platforms, then use \NewInfoField* with a star:
  % \NewInfoField*{mastodon}{\faMastodon}
  %% then you can use \mastodon, with TWO arguments where the 2nd argument is
  %% the full hyperlink.
  % \mastodon{@username@instance}{https://instance.url/@username}
}

\makecvheader

%% Depending on your tastes, you may want to make fonts of itemize environments slightly smaller
\AtBeginEnvironment{itemize}{\small}

%% Set the left/right column width ratio to 6:4.
\columnratio{0.6}

% Start a 2-column paracol. Both the left and right columns will automatically
% break across pages if things get too long.
\begin{paracol}{2}


\cvsection{Experience}

\cvevent{Graduate Researcher}{Human-Centered Computer Systems Lab, SNU}{Sept 2024 -- Present}{Seoul, South Korea}
\begin{itemize}
  \item Conducting research on \textbf{mobile AI privacy and accountable mobile systems}, focusing on how LLM/VLM-enabled apps infer and expose private attributes in real app workflows.
  \item Characterizing how model-derived attributes travel through \textbf{UI, storage, network, voice, and cross-app channels}, and developing system abstractions for \textbf{attribute-, workflow-, and channel-level privacy protection}.
  \item Earlier works include an \textbf{AI-driven scaffolding system for parent--child storytelling} and \textbf{on-demand affordance generation for cross-spatial object interaction}.
\end{itemize}

\divider

\cvevent{Teaching Assistant}{Seoul National University}{Sept 2025 -- Dec 2025}{Seoul, South Korea}
\begin{itemize}
  \item Served as a teaching assistant for \textbf{Software Development Principles and Practices (SWPP)}, supporting programming assignments, code reviews, and team-based software projects.
\end{itemize}

\divider

% \cvevent{Undergraduate Research Intern}{Human-Centered Computer Systems Lab, SNU}{Sept 2023 -- Aug 2024}{Seoul, South Korea}
% \begin{itemize}
%   \item Leveraged background in mathematics education to integrate social science research on child learning and scaffolding into system design.
%   \item Designed and built LLM-based systems for educational interaction, focusing on child–parent storytelling activities and narrative development.
%   \item Conducted qualitative and quantitative analysis of interaction data to evaluate engagement and learning outcomes.
% \end{itemize}

% \divider

\cvevent{Teaching Assistant}{Seoul National University}{Mar 2023 -- Jun 2023}{Seoul, South Korea}
\begin{itemize}
  \item Served as mentor and TA in a project-based course connecting \textbf{Education and Artificial Intelligence}. Guided student teams in designing AI-driven educational applications, providing feedback on both technical implementation and pedagogical grounding.
\end{itemize}



% \divider

% \cvevent{Product Engineer}{Google}{23 June 1999 -- 2001}{Palo Alto, CA}

% \begin{itemize}
% \item Joined the company as employe \#20 and female employee \#1
% \item Developed targeted advertisement in order to use user's search queries and show them related ads
% \end{itemize}

% \cvsection{A Day of My Life}

% % Adapted from @Jake's answer from http://tex.stackexchange.com/a/82729/226
% % \wheelchart{outer radius}{inner radius}{
% % comma-separated list of value/text width/color/detail}
% % Some ad-hoc tweaking to adjust the labels so that they don't overlap
% \hspace*{-1em}  %% quick hack to move the wheelchart a bit left
% \wheelchart{1.5cm}{0.5cm}{%
%   20/10em/accent!30/Running experiments \& debugging,
%   20/10em/accent!60/Discussing new ideas over coffee,
%   10/11em/accent!40/Organizing lab events,
%   10/11em/accent!20/Daydreaming about XR privacy futures,
%   10/9em/accent!50/Writing,
%   15/9em/accent!10/Reading papers,
%   5/10em/accent!40/Trying out new foods,
%   10/10em/accent!25/Social activities
% }

% % use ONLY \newpage if you want to force a page break for
% % ONLY the currentc column
% \newpage

\cvsection{Publications \& Presentations}

\cvpub
{TaleTrain: AI Scaffolding for Parent--Child Video Story Retelling at Home}
{Sangwon Park, \textbf{Sieun Park}, HyunA Seo, Minkyu Shim, Youngki Lee}
{Under Review}

\divider

\cvpub
{Stop the Hidden Inference: Privacy Safeguards Beyond Permissions}
{\textbf{Sieun Park}, Youngki Lee}
{Presented at the UIST 2025 Workshop}

\divider

\cvpub
{At Your Fingertips: On-Demand Affordances for Cross-Spatial Object Interaction}
{\textbf{Sieun Park}, Youngki Lee}
{Presented at CHI 2025 Workshop}

%% Switch to the right column. This will now automatically move to the second
%% page if the content is too long.
\switchcolumn

\cvsection{Life Philosophy}
\begin{quote}
``Research is not only about systems, but also about people — to support, connect, and inspire.''
\end{quote}

\cvsection{Most Proud of}

\cvachievement{\faLightbulb}{Research Vision}{for pioneering privacy-aware AI and XR systems that balance utility with user trust}

\divider

\cvachievement{\faUsers}{Community Building}{by fostering collaboration and communication in lab, volunteer, and academic communities}

\divider

\cvachievement{\faGlobe}{International Engagement}{through CHI/ UIST submissions, global volunteering, and presenting research to international audiences}

\cvsection{Strengths}

% Don't overuse these \cvtag boxes — they're just eye-candies and not essential. If something doesn't fit on a single line, it probably works better as part of an itemized list (probably inlined itemized list), or just as a comma-separated list of strengths.

\cvtag{Collaboration} 
\cvtag{Adaptability} 
\cvtag{Analytical Thinking}

\divider\smallskip

\cvtag{HCI \& XR Systems} 
\cvtag{AI Privacy} 
\cvtag{Interdisciplinary Approach}

\cvsection{Languages}

\cvskill{Korean}{5}
\divider

\cvskill{English}{4}



\cvsection{Education}

\cvevent{M.S./Ph.D\ in Artificial Intelligence}{Seoul National University}{Sept 2024 -- Present}{}

\divider

\cvevent{B.S.\ in Mathematics Education \& Artificial Intelligence}{Seoul National University}{Mar 2018 -- Feb 2024}{}

\newpage

% \cvsection{Referees}

% % \cvref{name}{email}{mailing address}
% \cvref{Prof.\ Alpha Beta}{Institute}{a.beta@university.edu}
% {Address Line 1\\Address Line 2}

% \divider

% \cvref{Prof.\ Gamma Delta}{Institute}{g.delta@university.edu}
% {Address Line 1\\Address Line 2}

\end{paracol}

\end{document}
